\chapter{Granica dolna i granica górna ciągu}
\label{ma1:ch:liminf-limsup}

Ciąg rzeczywisty nie musi mieć granicy. Może oscylować pomiędzy różnymi wartościami lub być nieograniczony. Granica dolna i granica górna pozwalają opisać jego zachowanie niezależnie od istnienia zwykłej granicy. Konstrukcja wykorzystuje porządek oraz kresy w rozszerzonej prostej rzeczywistej, wprowadzonej w rozdziale \ref{ma1:ch:real-extensions}. Nie przenosimy tych pojęć bezpośrednio na dowolną przestrzeń metryczną, w której może nie istnieć porządek.

\section{Zbieżność do nieskończoności}
\begin{definicja}
\label{ma1:def:extended-sequence-limit}
Niech $(a_n)_{n\ge1}$ będzie ciągiem w $\overline{\R}$. Mówimy, że $a_n$ zbiega do $+\infty$, jeżeli dla każdego $M\in\R$ istnieje $N\in\N_+$ takie, że dla wszystkich $n\ge N$ zachodzi $a_n>M$. Analogicznie $a_n\to-\infty$, jeżeli dla każdego $M\in\R$ istnieje $N\in\N_+$ takie, że dla wszystkich $n\ge N$ zachodzi $a_n<M$. Dla $a\in\R$ mówimy, że $a_n\to a$, jeżeli dla każdego rzeczywistego $\varepsilon>0$ od pewnego indeksu $a-\varepsilon<a_n<a+\varepsilon$.
\end{definicja}
W ostatnim przypadku wyrazy są od pewnego indeksu rzeczywiste, co wynika z warunku dla $\varepsilon=1$. Dla ciągów rzeczywistych definicja ta jest równoważna zbieżności w metryce wartości bezwzględnej. Zapis $\lim a_n=+\infty$ nie oznacza zbieżności w $\R$, lecz w sensie rozszerzonej prostej.

\begin{lemat}
\label{ma1:lem:extended-monotone}
Niech $(a_n)$ będzie ciągiem w $\overline{\R}$. Jeśli jest niemalejący, to zbiega w $\overline{\R}$ do $\sup_{n\ge1}a_n$. Jeśli jest nierosnący, to zbiega do $\inf_{n\ge1}a_n$.
\end{lemat}
\begin{proof}
Rozważmy ciąg niemalejący i oznaczmy $s=\sup_{n\ge1}a_n$, który istnieje na mocy twierdzenia \ref{ma1:thm:extended-complete}. Gdy $s=-\infty$, wszystkie wyrazy są równe $-\infty$. Gdy $s=+\infty$, dla każdego rzeczywistego $M$ istnieje indeks $N$ z $a_N>M$, ponieważ $M$ nie jest ograniczeniem górnym. Dla $n\ge N$ mamy $a_n\ge a_N>M$. Gdy $s\in\R$, liczba $s-\varepsilon$ nie jest ograniczeniem górnym dla żadnego $\varepsilon>0$. Wybierając $N$ z $a_N>s-\varepsilon$, otrzymujemy $s-\varepsilon<a_n\le s<s+\varepsilon$ dla $n\ge N$.
Dla ciągu nierosnącego stosujemy powyższy wynik do $(-a_n)$. Negacja odwraca porządek, zamienia kres dolny na przeciwny kres górny oraz zamienia zbieżność do $+\infty$ na zbieżność do $-\infty$ i odwrotnie. Dla skończonej granicy wynika to z nierówności definiujących zbieżność.
\end{proof}
\begin{przyklad}
Ciąg $a_n=n$ zbiega do $+\infty$ z własności Archimedesa. Ciąg $b_n=(-1)^nn$ jest nieograniczony zarówno z góry, jak i z dołu, lecz nie zbiega do żadnej z nieskończoności: ma dowolnie późne wyrazy dodatnie i ujemne. Nieograniczoność nie jest więc równoważna zbieżności do nieskończoności.
\end{przyklad}

\section{Kresy dalszych wyrazów}
\begin{definicja}
\label{ma1:def:liminf-limsup}
Niech $(a_n)_{n\ge1}$ będzie ciągiem w $\overline{\R}$. Dla $N\in\N_+$ oznaczamy
\[
 \alpha_N=\inf_{n\ge N}a_n,\qquad \beta_N=\sup_{n\ge N}a_n.
\]
Granicą dolną i granicą górną ciągu nazywamy odpowiednio
\[
 \liminf_{n\to\infty}a_n=\sup_{N\ge1}\alpha_N,\qquad
 \limsup_{n\to\infty}a_n=\inf_{N\ge1}\beta_N.
\]
Wszystkie kresy bierzemy w $\overline{\R}$.
\end{definicja}
\begin{twierdzenie}
\label{ma1:thm:tail-limits}
Dla dowolnego ciągu $(a_n)$ w $\overline{\R}$ ciąg $(\alpha_N)$ jest niemalejący, a $(\beta_N)$ nierosnący. Zachodzą
\[
 \liminf a_n=\lim_{N\to\infty}\alpha_N,\qquad
 \limsup a_n=\lim_{N\to\infty}\beta_N,\qquad
 \liminf a_n\le\limsup a_n.
\]
Granice po prawej stronie rozumiemy w $\overline{\R}$. Zmiana skończenie wielu początkowych wyrazów nie wpływa na żadną z tych wielkości.
\end{twierdzenie}
\begin{proof}
Zbiór $\{a_n:n\ge N+1\}$ zawiera się w $\{a_n:n\ge N\}$, zatem $\alpha_N\le\alpha_{N+1}$ i $\beta_{N+1}\le\beta_N$. Pierwsze dwie równości wynikają z lematu \ref{ma1:lem:extended-monotone}.
Dla dowolnych $N,M\ge1$ wybieramy $k=\max(N,M)$. Wtedy $\alpha_N\le a_k\le\beta_M$. Dla ustalonego $M$ liczba $\beta_M$ jest więc ograniczeniem górnym wszystkich $\alpha_N$, skąd $\sup_N\alpha_N\le\beta_M$. Biorąc kres dolny po $M$, dostajemy ostatnią nierówność.
Jeśli dwa ciągi mają wspólne wyrazy od indeksu $N_0$, ich kresy dalszych wyrazów są równe dla $N\ge N_0$. Granice ciągów tych kresów są więc równe.
\end{proof}

\begin{przyklad}
\label{ma1:ex:oscillating-limits}
Niech $a_{2k-1}=1+1/k$ i $a_{2k}=-1-1/k$ dla $k\ge1$. Dla $N=2k-1$ mamy $\alpha_N=-1-1/k$ i $\beta_N=1+1/k$. Istotnie, w dalszych wyrazach dodatnia część maleje, a ujemna rośnie. Dla $N=2k$ mamy $\alpha_N=-1-1/k$, $\beta_N=1+1/(k+1)$. Zatem $\liminf a_n=-1$ i $\limsup a_n=1$.
\end{przyklad}

\begin{figure}[htbp]
\centering
\begin{tikzpicture}[x=.64cm,y=.85cm,>=Stealth]
\draw[->] (0,0)--(13.4,0) node[right] {$n$};
\draw[->] (0,-2.65)--(0,2.65) node[above] {$a_n$};
\draw[dashed,gray] (0,1)--(12.7,1);
\draw[dashed,gray] (0,-1)--(12.7,-1);
\node[left] at (0,1) {$1$};
\node[left] at (0,-1) {$-1$};
\foreach \k in {1,...,6}{
 \pgfmathsetmacro{\u}{1+1/\k}
 \pgfmathsetmacro{\v}{-1-1/\k}
 \pgfmathtruncatemacro{\odd}{2*\k-1}
 \pgfmathtruncatemacro{\even}{2*\k}
 \fill[blue!70!black] (\odd,\u) circle (2pt);
 \fill[blue!70!black] (\even,\v) circle (2pt);
}
\draw[orange!85!black,thick] (1,2)--(2,2)--(2,1.5)--(4,1.5)--(4,{1+1/3})--(6,{1+1/3})--(6,1.25)--(8,1.25)--(8,1.2)--(10,1.2)--(10,{1+1/6})--(12,{1+1/6});
\draw[teal!75!black,thick] (1,-2)--(3,-2)--(3,-1.5)--(5,-1.5)--(5,{-1-1/3})--(7,{-1-1/3})--(7,-1.25)--(9,-1.25)--(9,-1.2)--(11,-1.2)--(11,{-1-1/6})--(12,{-1-1/6});
\node[orange!85!black,above] at (7,1.5) {$\beta_N$};
\node[teal!75!black,below] at (7,-1.5) {$\alpha_N$};
\node[below left] at (0,0) {$0$};
\end{tikzpicture}
\caption{Wyrazy ciągu z przykładu \ref{ma1:ex:oscillating-limits} oraz schodkowa ilustracja kresów jego dalszych wyrazów.}
\label{ma1:fig:tail-limits}
\end{figure}
\FloatBarrier

\section{Kryteria za pomocą progów}
\begin{twierdzenie}
\label{ma1:thm:limsup-thresholds}
Niech $(a_n)$ będzie ciągiem w $\overline{\R}$, a $L=\limsup a_n$. Dla dowolnego $t\in\R$ zachodzą:
\begin{enumerate}[label=(\roman*)]
\item jeśli $t>L$, to od pewnego indeksu $a_n<t$;
\item jeśli $t<L$, to dla każdego $N\ge1$ istnieje $n\ge N$ z $a_n>t$.
\end{enumerate}
Te dwie własności charakteryzują $L$ jednoznacznie w $\overline{\R}$. Dla $l=\liminf a_n$ zachodzą odpowiednio: jeśli $t<l$, to od pewnego indeksu $a_n>t$; jeśli $t>l$, to dla każdego $N$ istnieje $n\ge N$ z $a_n<t$.
\end{twierdzenie}
\begin{proof}
Jeśli $t>L=\inf_N\beta_N$, to $t$ nie jest ograniczeniem dolnym zbioru $\{\beta_N:N\ge1\}$. Istnieje $N$ z $\beta_N<t$, a wtedy $a_n\le\beta_N<t$ dla $n\ge N$.
Jeśli $t<L$, to dla każdego $N$ mamy $\beta_N\ge L>t$. Gdyby wszystkie $a_n$ dla $n\ge N$ były nie większe od $t$, to $\beta_N\le t$, sprzeczność.
Dla jednoznaczności przypuśćmy, że $L_1<L_2$ spełniają obie własności. Istnieje rzeczywiste $t$ z $L_1<t<L_2$: dla skończonych końców bierzemy środek, dla jednego nieskończonego końca przesuwamy skończony koniec o $1$, a dla dwóch nieskończonych końców bierzemy $0$. Z własności (i) dla $L_1$ od pewnego indeksu $a_n<t$, podczas gdy (ii) dla $L_2$ zapewnia dowolnie późne wyrazy większe od $t$. Sprzeczność. Dla granicy dolnej stosujemy tezy do $(-a_n)$ i korzystamy z odwracania porządku.
\end{proof}
W szczególności dla $L\in\R$ i każdego $\varepsilon>0$ od pewnego indeksu $a_n<L+\varepsilon$, a dowolnie późno zachodzi $a_n>L-\varepsilon$. Nie twierdzimy, że dowolnie późno $a_n=L$ ani że od pewnego indeksu $a_n\le L$.

\section{Związek z podciągami}
\begin{twierdzenie}
\label{ma1:thm:extremal-subsequence-limits}
Niech $(a_n)$ będzie ciągiem w $\overline{\R}$ oraz $l=\liminf a_n$, $L=\limsup a_n$. Każda granica podciągu w $\overline{\R}$ należy do przedziału $[l,L]$. Ponadto istnieją podciągi zbieżne odpowiednio do $l$ i do $L$. W szczególności $l$ i $L$ są najmniejszą i największą granicą podciągu w $\overline{\R}$.
\end{twierdzenie}
\begin{proof}
Jeśli $a_{n_k}\to b$ i $b>L$, wybieramy rzeczywiste $t$ pomiędzy $L$ i $b$. Z twierdzenia \ref{ma1:thm:limsup-thresholds} od pewnego indeksu $a_n<t$, a z definicji zbieżności do $b$ od pewnego indeksu $k$ zachodzi $a_{n_k}>t$. Dla $b\in\R$ używamy dokładności $(b-t)/2$, dla $b=+\infty$ bezpośrednio progu $t$. Ponieważ $n_k\ge k$, otrzymujemy sprzeczność. Argument z negacją wyklucza $b<l$.

Załóżmy najpierw $L\in\R$. Dla każdego $k\ge1$ istnieje $M_k$ takie, że $a_n<L+1/k$ dla wszystkich $n\ge M_k$. Przyjmijmy $n_0=0$. Z drugiej własności progowej wybieramy $n_k\ge\max(M_k,n_{k-1}+1)$ z $a_{n_k}>L-1/k$. Wówczas $L-1/k<a_{n_k}<L+1/k$, zatem podciąg zbiega do $L$.
Jeśli $L=+\infty$, dla każdego $k$ wybieramy $n_k>n_{k-1}$ z $a_{n_k}>k$. Dla dowolnego rzeczywistego $M$ istnieje $K$ z $K>M$, więc dla $k\ge K$ mamy $a_{n_k}>M$. Jeśli $L=-\infty$, z pierwszej własności progowej dla każdego $t\in\R$ od pewnego indeksu $a_n<t$. W tym przypadku cały ciąg zbiega do $-\infty$ i sam jest odpowiednim podciągiem. Konstrukcja podciągu zbieżnego do $l$ wynika z zastosowania powyższych argumentów do $(-a_n)$.
\end{proof}
\begin{przyklad}
Ciąg $(-1)^n$ ma granicę dolną $-1$ i górną $1$, ale żadna liczba z otwartego przedziału $(-1,1)$ nie jest granicą jego podciągu. Twierdzenie mówi o krańcach zbioru granic podciągów, a nie o wypełnieniu całego przedziału.
\end{przyklad}

\begin{twierdzenie}
\label{ma1:thm:liminf-limsup-convergence}
Niech $(a_n)$ będzie ciągiem w $\overline{\R}$ i $a\in\overline{\R}$. Wówczas $a_n\to a$ wtedy i tylko wtedy, gdy $\liminf a_n=\limsup a_n=a$.
\end{twierdzenie}
\begin{proof}
Jeżeli $a_n\to a$ i $a\in\R$, to dla $\varepsilon>0$ od pewnego indeksu $N$ zachodzi $a-\varepsilon<a_n<a+\varepsilon$. Stąd $\alpha_N\ge a-\varepsilon$, $\beta_N\le a+\varepsilon$, czyli $a-\varepsilon\le l\le L\le a+\varepsilon$. Ponieważ zachodzi to dla każdego $\varepsilon>0$, mamy $l=L=a$. Jeśli $a=+\infty$, dla każdego $M\in\R$ od pewnego indeksu $a_n>M$, zatem $\alpha_N\ge M$ i $l\ge M$. Wobec tego $l=+\infty$ i $L\ge l$ daje $L=+\infty$. Przypadek $a=-\infty$ wynika przez negację.
Odwrotnie, jeśli $l=L=a\in\R$, stosujemy kryteria progowe do $a-\varepsilon$ oraz $a+\varepsilon$. Od wspólnego indeksu zachodzi $a-\varepsilon<a_n<a+\varepsilon$, co daje zbieżność. Dla $l=L=+\infty$ kryterium granicy dolnej daje dla każdego $M$ nierówność $a_n>M$ od pewnego indeksu; dla $-\infty$ korzystamy z kryterium granicy górnej.
\end{proof}
